% rhythm: article display math opening a paragraph: a display after a list end (no empty line, TeX's \@endpe state), a display after a \vspace and a display after another display (the empty line amsmath's $$ sets in vertical mode) — the facts the engine reads from the source, not from its owed glue
% boundary: listend-display Cedar span
% boundary: vspace-display Fir span
% boundary: display-display Ivy span
\documentclass{article}
\usepackage{amsmath}
\usepackage{fontspec}
\setmainfont{OpenSans-Regular.ttf}[Path=../corpus/fonts/, BoldFont=OpenSans-Bold.ttf, ItalicFont=OpenSans-Italic.ttf, BoldItalicFont=OpenSans-BoldItalic.ttf]
\usepackage{unicode-math}
\setmathfont{FiraMath-Regular.otf}[Path=../corpus/fonts/]
\begin{document}
Alder sets a line before the list.
\begin{itemize}
\item Birch item of the list.
\end{itemize}
\[ c + d = \text{Cedar} \]
Dogwood follows the display after the list end.

Elm sets one line before the vspace.

\vspace{12pt}

\[ e + f = \text{Fir} \]
Grove follows the display after the vspace.

Hazel sets one line before the pair of displays.

\[ g + h = \text{Holly} \]

\[ i + j = \text{Ivy} \]
Juniper follows the second display.
\end{document}
